\section{Stability: the Gram matrix survives small perturbations}\label{sec:gram-stability}

\begin{lemma}[Gram stability under weight perturbation]\label{lem:gram-stability}
Fix $\delta \in (0,1)$ and a radius
\begin{align}\label{eq:radius}
R \;=\; \frac{c\, \delta\, \lzero}{n^2},
\end{align}
for a sufficiently small universal constant $c > 0$. Let $W = (w_1, \dots, w_m)$ be any
weights (possibly depending on the data) with $\norm{w_r - w_r(0)} \le R$ for every
$r \in [m]$. Then with probability at least $1 - \delta$ over the random initialization
$W(0)$,
\begin{align}\label{eq:gram-stable}
\opnorm{H(W) - H(0)} \;\le\; \frac{\lzero}{4}.
\end{align}
Consequently, on the intersection with the event of \Cref{lem:init-gram-close},
$\lmin(H(W)) \ge \tfrac{1}{2}\lzero$.
\end{lemma}

\begin{proof}
The strategy is: identify the event that lets a neuron change the Gram matrix (a sign flip of
its activation), bound its probability by Gaussian anti-concentration, take expectation to
control $\E\opnorm{H(W)-H(0)}$, and convert to high probability by Markov.

\smallskip\noindent\textbf{Step 1 (a neuron contributes only if it flips).}
With $\sigma$ the ReLU, $\indic{w_r^{\top} x_i \ge 0}$ is the activation indicator of neuron
$r$ on $x_i$. Subtracting Eq.~\eqref{eq:emp-gram} at $W$ and at $W(0)$,
\begin{align}\label{eq:entry-diff}
\abs{H_{ij}(W) - H_{ij}(0)}
\;\le\; & ~ \frac{1}{m} \sum_{r=1}^{m} \abs{x_i^{\top} x_j} \cdot
\abs{ S_{r,ij}(W) - S_{r,ij}(0) } \\
\;\le\; & ~ \frac{1}{m} \sum_{r=1}^{m} \indic{ E_{r,i} \cup E_{r,j} },
\end{align}
where $S_{r,ij}(\cdot) := \indic{w_r^{\top} x_i \ge 0}\indic{w_r^{\top} x_j \ge 0}$, the first
step is the triangle inequality applied termwise, and the second uses $\abs{x_i^{\top}x_j}\le1$
together with the observation that $S_{r,ij}(W) \ne S_{r,ij}(0)$ forces neuron $r$ to flip its
sign on $x_i$ or on $x_j$; here
\begin{align}\label{eq:flip-event}
E_{r,i} \;:=\; \left\{ \abs{ w_r(0)^{\top} x_i } \le R \right\}.
\end{align}
The inclusion holds because, by Cauchy--Schwarz and $\norm{x_i}=1$,
$\abs{(w_r - w_r(0))^{\top} x_i} \le \norm{w_r - w_r(0)} \le R$, so a sign change of
$w_r^{\top} x_i$ relative to $w_r(0)^{\top} x_i$ requires $\abs{w_r(0)^{\top} x_i} \le R$.

\smallskip\noindent\textbf{Step 2 (Gaussian anti-concentration).}
Since $\norm{x_i} = 1$, the projection $w_r(0)^{\top} x_i \sim \N(0, 1)$. The standard normal
density is bounded by $1/\sqrt{2\pi} \le 1/2$, so for the small-ball event Eq.~\eqref{eq:flip-event},
\begin{align}\label{eq:anti-conc}
\Pr\!\left[ E_{r,i} \right]
\;=\; & ~ \Pr\!\left[ \abs{\N(0,1)} \le R \right]
\;=\; \int_{-R}^{R} \frac{1}{\sqrt{2\pi}}\, e^{-t^2/2}\, dt
\;\le\; \frac{2R}{\sqrt{2\pi}}
\;\le\; R,
\end{align}
where the first equality is the definition of $E_{r,i}$, the second writes the Gaussian
probability as an integral, the third bounds the integrand by its maximum $1/\sqrt{2\pi}$ over
the interval of length $2R$, and the last uses $2/\sqrt{2\pi} \le 1$.

\smallskip\noindent\textbf{Step 3 (a $W$-independent dominating quantity).}
The right-hand side of Eq.~\eqref{eq:entry-diff} depends on $W(0)$ alone, not on $W$, and
dominates $\abs{H_{ij}(W)-H_{ij}(0)}$ for \emph{every} $W$ in the ball simultaneously. Define
the single random variable
\begin{align}\label{eq:dom-quantity}
D \;:=\; \sum_{i=1}^{n}\sum_{j=1}^{n} \frac{1}{m}\sum_{r=1}^{m}
\indic{ E_{r,i} \cup E_{r,j} },
\end{align}
which is a function of $W(0)$ only. By Eq.~\eqref{eq:entry-diff} and
$\opnorm{\cdot} \le \fnorm{\cdot} \le \sum_{ij}\abs{\cdot}$,
\begin{align}\label{eq:dom-sup}
\sup_{\,W :\; \norm{w_r-w_r(0)}\le R\;\forall r}\; \opnorm{H(W)-H(0)} \;\le\; D .
\end{align}
Taking expectations
and using the union bound $\Pr[E_{r,i}\cup E_{r,j}] \le \Pr[E_{r,i}] + \Pr[E_{r,j}] \le 2R$
together with Eq.~\eqref{eq:anti-conc} for each of the $m$ neurons,
\begin{align}\label{eq:exp-op}
\E[D]
\;=\; & ~ \sum_{i=1}^{n}\sum_{j=1}^{n} \frac{1}{m}\sum_{r=1}^{m}
\Pr\!\left[ E_{r,i} \cup E_{r,j} \right]
\;\le\; \sum_{i=1}^{n}\sum_{j=1}^{n} 2R
\;=\; 2 n^2 R,
\end{align}
where the first step is linearity of expectation, the second applies Eq.~\eqref{eq:anti-conc}
to each neuron, and the last counts the $n^2$ entries.

\smallskip\noindent\textbf{Step 4 (Markov to high probability, uniformly over $W$).}
By Markov's inequality applied to the nonnegative $W$-independent random variable $D$,
\begin{align}\label{eq:markov-stab}
\Pr\!\left[ D \ge \frac{\lzero}{4} \right]
\;\le\; & ~ \frac{4}{\lzero}\, \E[D]
\;\le\; \frac{8 n^2 R}{\lzero}
\;=\; 8 c\, \delta
\;\le\; \delta,
\end{align}
where the first step is Markov, the second substitutes Eq.~\eqref{eq:exp-op}, the third
substitutes the radius Eq.~\eqref{eq:radius}, and the last holds for $c \le 1/8$. On the
complementary event $\{D \le \lzero/4\}$, which has probability at least $1-\delta$ and depends
on $W(0)$ alone, $\opnorm{H(W)-H(0)} \le D \le \lzero/4$ holds for \emph{every} admissible $W$
at once; in particular it applies to the data-dependent iterates $W(k)$ downstream. This
establishes Eq.~\eqref{eq:gram-stable} as the uniform event claimed.

\smallskip\noindent\textbf{Step 5 (eigenvalue floor).}
On the intersection with the event of \Cref{lem:init-gram-close} (where
$\lmin(H(0)) \ge \tfrac34\lzero$), Weyl's inequality (\Cref{fac:weyl}) gives
\begin{align}\label{eq:weyl-stab}
\lmin\!\left( H(W) \right)
\;\ge\; & ~ \lmin\!\left( H(0) \right) - \opnorm{H(W) - H(0)}
\;\ge\; \frac{3}{4}\lzero - \frac{1}{4}\lzero
\;=\; \frac{1}{2}\lzero,
\end{align}
where the first step is Weyl applied to the symmetric matrices $H(W)$ and $H(0)$ and the
second substitutes the two events. This gives the assertion.
\end{proof}
